\documentclass[10pt,a4paper]{amsart}
\usepackage[margin=19mm]{geometry}
\usepackage{amsmath,amssymb,mathtools}
\usepackage[scr=rsfs,cal=euler]{mathalfa}
\usepackage{xfrac}
\usepackage[unicode,hidelinks,bookmarks=false]{hyperref}
\newcommand{\ie}{i.e.\ }
\newcommand{\cf}{cf.\ }
\newcommand{\resp}{respectively\ }
\newcommand{\Subsection}{\subsection}
\providecommand{\nobreakdash}{\nobreak\hbox{-}}
\setlength{\emergencystretch}{2em}
\multlinegap0pt
\usepackage{amssymb}
\usepackage{mathtools}
\usepackage[shortlabels]{enumitem}
\newtheorem{lemma}{Lemma}
\newcommand{\cA}{\mathcal A}
\title{Simplicial arrangements in real projective three-space revisited}
\author{Marek Janasz, Piotr Pokora}
\date{Source: arXiv:2608.28254v1 (CC BY 4.0)}
\begin{document}
\maketitle

\noindent\emph{Source and licence.} Simplicial arrangements in real projective three-space revisited. Version arXiv:2608.28254v1 (CC BY 4.0), distributed by arXiv under the Creative Commons Attribution 4.0 International licence, http://creativecommons.org/licenses/by/4.0/, as recorded in the arXiv licence metadata for this exact version and shown on the abstract page https://arxiv.org/abs/2608.28254v1. Full licence notice: \texttt{licenses/arxiv-2608.28254-CC-BY-4.0.txt}.\par\smallskip

\noindent\textit{Editorial scope.} Counted unit: the article's lemma lem:face-numbers-incidence (source0.tex lines 619--642). The article's complete original proof of it is delivered with it (source0.tex lines 644--705, one \texttt{proof} environment of 62 lines). No mathematics is written, repaired or re-derived: the statement and the proof are the article's own lines, in the article's own notation, and no retained line is substituted or re-flowed.  No further source-local context is needed: every cross-reference of every retained line resolves inside the two delivered spans.  Nothing was omitted from any retained span, so the package carries no omission record.  No retained line contains a citation, so the wrapper carries no bibliography and no bibliography entry.  No retained line contains a figure environment, an image-inclusion command or drawing code, so no image is delivered and the package declares no asset dependency.  Editorial formatting only, in addition: an AMS \texttt{amsart} preamble that restates the article's own declarations and macros used by the retained lines, namely the environments \texttt{lemma} on separate counters, together with the standard packages those lines need.  The retained environments print in this document as Lemma 1 in order of appearance, whereas the article prints the same items with the numbering of its own full document.  The complete native source of the article as distributed in the arXiv e-print for this exact version is bundled unchanged as \texttt{sources/208774/source0.tex}.
\begin{lemma}
\label{lem:face-numbers-incidence}
With the notation above, one has
\[
f_0
=
\sum_{p\geq 3}t_p,
\]
\[
f_1
=
\sum_{p\geq 3}\sum_{q\geq 2}t_{pq},
\]
and
\[
f_2
=
n
+
\sum_{p\geq 3}\sum_{q\geq 2}q\,t_{pq}
-
\sum_{p\geq 3}p\,t_p.
\]
\end{lemma}

\begin{proof}
The first identity is immediate. We first note that every
$L\in L_2(\cA)$ contains a vertex. Indeed, in the central picture let
$X$ be the two-dimensional subspace corresponding to $L$. The
restrictions to $X$ of the defining forms of the hyperplanes not
containing $X$ must span $X^*$; otherwise they would have a nonzero
common kernel in $X$, and that vector would lie in every hyperplane of
the arrangement, contradicting essentiality. Hence at least two
distinct rank-three flats lie on $L$.

For the second identity, every one-dimensional cell lies on a unique
intersection line. A projective line containing $r$ vertices is divided
by them into exactly $r$ one-dimensional cells.
Consequently, $f_1$ is the total number of incidences between vertices
and intersection lines, which is
\[
\sum_{p\geq 3}\sum_{q\geq 2}t_{pq}.
\]
For $H\in\cA$, let $\cA^H$ be the reduced restriction and, for a vertex
$x\in H$, put
\[
\lambda_H(x)
=
\bigl|\{L\in L_2(\cA)\mid x\in L\subset H\}\bigr|.
\]
Again the Euler formula for the projective line arrangement $\cA^H$ gives us
\[
r(\cA^H)
=
1+
\sum_{\substack{x\in L_3(\cA)\\x\in H}}
\bigl(\lambda_H(x)-1\bigr).
\]
Summing over $H\in\cA$ and using the fact
\[
f_2=\sum_{H\in\cA}r(\cA^H),
\]
we obtain
\[
f_2
=
n+
\sum_{x\in L_3(\cA)}
\left(
\sum_{\substack{L\in L_2(\cA)\\x\in L}}m(L)-\nu(x)
\right).
\]
Finally,
\[
\sum_{x\in L_3(\cA)}\nu(x)
=
\sum_{p\geq 3}p\,t_p
\]
and
\[
\sum_{x\in L_3(\cA)}
\sum_{\substack{L\in L_2(\cA)\\x\in L}}m(L)
=
\sum_{p\geq 3}\sum_{q\geq 2}q\,t_{pq},
\]
which gives the asserted formula for $f_2$.
\end{proof}

\end{document}
